$\NestedShrinkIntervalUp$ consumes $\CauchyCompletionCriterionUp$ after the shrinking-diameter ledger has been repacked through the $\StreamNameUp$ observation schedule, the $\RegSeqRatUp$ regular rational readback, the $\DyadicRatCoreUp$ tolerance ledger, and the terminal $\RealUp$ seal route. This sibling dependency records the completion-criterion consumer path of \autoref{ch:concrete-instances-cauchycompletioncriterion-namecert}: a nested shrink interval can be read as completion evidence only after its finite diameter data has passed through the same stream, regular sequence, dyadic tolerance, and Real seal rows.

\begin{theorem}[Nested shrink interval Cauchy completion consumer]
\label{thm:nested-shrink-interval-cauchy-completion-consumer}
Every completion-facing read of $\NestedShrinkIntervalUp$ consumes $\CauchyCompletionCriterionUp$ only after the $\StreamNameUp$ observation schedule, the $\RegSeqRatUp$ regular rational readback, the $\DyadicRatCoreUp$ tolerance ledger, and the terminal $\RealUp$ seal rows are visible. Thus the dependency on \autoref{ch:concrete-instances-cauchycompletioncriterion-namecert} is a theorem-level consumer edge through \autoref{ch:concrete-instances-streamname-namecert} and \autoref{ch:concrete-instances-regseqrat-namecert}, not an appeal to an undisplayed completion object.
\end{theorem}
\begin{proof}
Read a completion-facing request for a nested shrink interval. The route first displays the finite observation schedule and regular rational readback rows, so the consumer sees only finite windows and readback entries.

The dyadic tolerance ledger and terminal Real seal are then exposed before the completion criterion is consumed. Since the criterion is reached after those rows, the dependency records a reusable finite route rather than a hidden selected limit or ambient completeness principle.
\end{proof}
